\chapter{Fermionic operators and signs}
\label{chap:operators}

\section{One global ordering}

Fermionic signs are meaningless until the spin orbitals have a global order.
This project fixes
\begin{equation}
(0\uparrow,1\uparrow,\ldots,(L-1)\uparrow,
 0\downarrow,1\downarrow,\ldots,(L-1)\downarrow).
\label{eq:ordering}
\end{equation}
Thus site $i$ has global orbital index $i$ for up spin and $L+i$ for down
spin.  The two state integers can be viewed as one $2L$-bit integer,
\begin{equation}
B=b_\uparrow+(b_\downarrow\mathbin{\texttt{<<}}L).
\end{equation}

For orbital $p$, define the parity count
\begin{equation}
m_p(B)=\operatorname{popcount}\left(B\mathbin{\texttt{\&}}(2^p-1)\right).
\end{equation}
It counts occupied orbitals strictly before $p$.  The elementary actions are
\begin{align}
c_p\ket{B} &=( -1)^{m_p(B)}n_p\ket{B\text{ with }p\text{ cleared}},\\
c_p^\dagger\ket{B} &=( -1)^{m_p(B)}(1-n_p)
\ket{B\text{ with }p\text{ set}}.
\end{align}
An invalid annihilation or creation returns no state rather than a zero-amplitude
state.

\section{A hop is two operator actions}

The hopping operator $c_p^\dagger c_q$ acts from right to left:
\begin{enumerate}
\item annihilate at $q$ using the original occupation;
\item create at $p$ using the intermediate occupation;
\item multiply the two signs.
\end{enumerate}
Computing a sign from a shortcut tied only to geometric distance is fragile.
The two-action algorithm remains correct for different spins, long-range
hopping, and periodic boundary bonds.

Consider $L=4$, $b_\uparrow=1010_2$, and the periodic hop $3\to0$.  Site $1$
is occupied between orbitals $0$ and $3$, so
\begin{equation}
c_{0\uparrow}^\dagger c_{3\uparrow}ket{1010_2}=-\ket{0011_2}.
\end{equation}
The sign is not $+1$.  This is the boundary case most likely to expose an
incorrect nearest-neighbor shortcut.

\lstinputlisting[caption={Inspecting creation and periodic hopping signs.},label={lst:signs}]{code/step02_fermion_signs.py}

\section{Algebraic tests}

The implementation is checked against the canonical anticommutators
\begin{align}
\{c_p,c_q^\dagger\}&=\delta_{pq},\\
\{c_p,c_q\}&=0,\\
\{c_p^\dagger,c_q^\dagger\}&=0.
\end{align}
For small $L$, these identities can be tested exhaustively on the complete
$2^{2L}$-state Fock space.  Exhaustive tests are inexpensive here and stronger
than testing only a few hand-selected signs.

The down-spin sign includes all occupied up orbitals because the down block
comes second in Eq.~\eqref{eq:ordering}.  In a number-conserving down-spin hop,
the up-block contribution appears twice and cancels, but creation and
annihilation must still each obey the global convention.  This distinction
becomes important as soon as non-number-conserving operators or Green
functions are introduced.

\paragraph{Checkpoint.}
Do not construct a Hamiltonian until the full small-space anticommutators,
cross-spin signs, an adjacent hop, and an occupied periodic-boundary hop all
pass.
